\documentclass[10pt,a4paper]{amsart}
\usepackage{amsmath,amssymb,amsthm,graphicx}
% Retained preamble for the Stacks Project chapter divisors.tex (Divisors).
% The chapter file pinned-original-source.tex carries only \input{preamble} on its first line, so no
% preamble text lives inside the chapter itself.  The Stacks Project uses one single shared preamble for
% every chapter; the snapshot stages it verbatim in sources/stacks-chow/preamble.tex.  The block below is
% copied mechanically, line for line, from that file's lines 14-18 (the Xy-pic setup) and lines 154-217
% (the native theorem declarations together with the mathematical macros).  Nothing was added, removed or
% redefined.  The wrapper class and the amsmath/amssymb/amsthm/graphicx base packages are supplied by the
% assembler, so the class/\externaldocument part of the shared preamble is not repeated here.
% For commutative diagrams we use Xy-pic
\usepackage[all]{xy}

% We use 2cell for 2-commutative diagrams.
\xyoption{2cell}
% Theorem environments.
%
\theoremstyle{plain}
\newtheorem{theorem}[subsection]{Theorem}
\newtheorem{proposition}[subsection]{Proposition}
\newtheorem{lemma}[subsection]{Lemma}

\theoremstyle{definition}
\newtheorem{definition}[subsection]{Definition}
\newtheorem{example}[subsection]{Example}
\newtheorem{exercise}[subsection]{Exercise}
\newtheorem{situation}[subsection]{Situation}

\theoremstyle{remark}
\newtheorem{remark}[subsection]{Remark}
\newtheorem{remarks}[subsection]{Remarks}

\numberwithin{equation}{subsection}

% Macros
%
\def\lim{\mathop{\mathrm{lim}}\nolimits}
\def\colim{\mathop{\mathrm{colim}}\nolimits}
\def\Spec{\mathop{\mathrm{Spec}}}
\def\Hom{\mathop{\mathrm{Hom}}\nolimits}
\def\Ext{\mathop{\mathrm{Ext}}\nolimits}
\def\SheafHom{\mathop{\mathcal{H}\!\mathit{om}}\nolimits}
\def\SheafExt{\mathop{\mathcal{E}\!\mathit{xt}}\nolimits}
\def\Sch{\mathit{Sch}}
\def\Mor{\mathop{\mathrm{Mor}}\nolimits}
\def\Ob{\mathop{\mathrm{Ob}}\nolimits}
\def\Sh{\mathop{\mathit{Sh}}\nolimits}
\def\NL{\mathop{N\!L}\nolimits}
\def\CH{\mathop{\mathrm{CH}}\nolimits}
\def\proetale{{pro\text{-}\acute{e}tale}}
\def\etale{{\acute{e}tale}}
\def\QCoh{\mathit{QCoh}}
\def\Ker{\mathop{\mathrm{Ker}}}
\def\Im{\mathop{\mathrm{Im}}}
\def\Coker{\mathop{\mathrm{Coker}}}
\def\Coim{\mathop{\mathrm{Coim}}}

% Boxtimes
%
\DeclareMathSymbol{\boxtimes}{\mathbin}{AMSa}{"02}

%
% Macros for moduli stacks/spaces
%
\def\QCohstack{\mathcal{QC}\!\mathit{oh}}
\def\Cohstack{\mathcal{C}\!\mathit{oh}}
\def\Spacesstack{\mathcal{S}\!\mathit{paces}}
\def\Quotfunctor{\mathrm{Quot}}
\def\Hilbfunctor{\mathrm{Hilb}}
\def\Curvesstack{\mathcal{C}\!\mathit{urves}}
\def\Polarizedstack{\mathcal{P}\!\mathit{olarized}}
\def\Complexesstack{\mathcal{C}\!\mathit{omplexes}}
% \Pic is the operator that assigns to X its picard group, usage \Pic(X)
% \Picardstack_{X/B} denotes the Picard stack of X over B
% \Picardfunctor_{X/B} denotes the Picard functor of X over B
\def\Pic{\mathop{\mathrm{Pic}}\nolimits}
\def\Picardstack{\mathcal{P}\!\mathit{ic}}
\def\Picardfunctor{\mathrm{Pic}}
\def\Deformationcategory{\mathcal{D}\!\mathit{ef}}

% --- appended 2026-10-12 (Chief): the project own definitions of the reference/slogan/history
% display environments, copied mechanically from the shared Stacks preamble lines 8-12
% (\usepackage{verbatim} plus the three \newenvironment lines, which the chapter preamble keeps
% verbatim from the project preamble). Nothing added, removed or redefined: these environments
% carry no mathematical content and the project itself typesets them as comments. Without these
% three lines a counted statement whose own text uses them fails with Environment slogan undefined.
\usepackage{verbatim}
\newenvironment{reference}{\comment}{\endcomment}
\newenvironment{slogan}{\comment}{\endcomment}
\newenvironment{history}{\comment}{\endcomment}
\usepackage{cleveref}
\crefname{originallemma}{original lemma}{original lemmas}
\Crefname{originallemma}{Original lemma}{Original lemmas}
\makeatletter
\newlabel{lemma-norm-map-invertible}{{0BCZ}{source}{}{sourcepointer-1}{}}
\newlabel{lemma-norm-map-invertible@cref}{{[originallemma][0][]0BCZ}{[1][0][]source}{}{}{}}
\makeatother
\crefname{originaldefinition}{original definition}{original definitions}
\Crefname{originaldefinition}{Original definition}{Original definitions}
\makeatletter
\newlabel{properties-definition-ample}{{01PS}{source}{}{sourcepointer-2}{}}
\newlabel{properties-definition-ample@cref}{{[originaldefinition][0][]01PS}{[1][0][]source}{}{}{}}
\makeatother
\crefname{originallemma}{original lemma}{original lemmas}
\Crefname{originallemma}{Original lemma}{Original lemmas}
\makeatletter
\newlabel{properties-lemma-affine-cap-s-open}{{01PV}{source}{}{sourcepointer-3}{}}
\newlabel{properties-lemma-affine-cap-s-open@cref}{{[originallemma][0][]01PV}{[1][0][]source}{}{}{}}
\makeatother
\crefname{originallemma}{original lemma}{original lemmas}
\Crefname{originallemma}{Original lemma}{Original lemmas}
\makeatletter
\newlabel{properties-lemma-ample-finite-set-in-principal-affine}{{09NV}{source}{}{sourcepointer-4}{}}
\newlabel{properties-lemma-ample-finite-set-in-principal-affine@cref}{{[originallemma][0][]09NV}{[1][0][]source}{}{}{}}
\makeatother
\title{The norm of an ample invertible sheaf along a finite morphism with a degree d norm is ample on the base}
\author{The Stacks Project Authors}
\date{stacks-project:divisors.tex:local-snapshot-20261012; GNU FDL 1.2 or later}
\begin{document}
\maketitle
\noindent\textit{Editorial note.} All delivered mathematical text is the author's own, copied from the pinned Stacks Project chapter file pinned-original-source.tex; no statement or proof step was rewritten, reordered or removed.\par\medskip
\noindent\textit{Editorial note.} Following the multi-claim convention for this source, source\_url is the source-level URL https://stacks.math.columbia.edu/ and the source label lemma-norm-ample is carried in claim\_key and here rather than in a per-claim tag URL.\par\medskip
\noindent\textit{Editorial note.} The preamble is the chapter's shared Stacks Project preamble, retained as sources/stacks-divisors/retained-preamble.tex: a line-for-line excerpt of the shared Stacks preamble staged with this snapshot. No macro was invented or redefined, and this card adds no preamble text of its own.\par\medskip
\noindent\textit{Editorial note.} Statement span 3737-3742 and proof span 3744-3770 come from the halfA pool card operations/pool-cards/stacks-divisors-halfA.json and were re-read in the pinned chapter file before this card was written.\par\medskip
\noindent\textit{Editorial note.} The proof names the norm sheaf and uses the author's notation for the norm of a section; both are preserved unchanged.\par\medskip
\noindent\textit{Editorial note.} The proof's conclusion uses the cap criterion with the phrase that the nonvanishing locus is affine; the author's wording is kept.\par\medskip
\noindent\textit{Editorial note.} 4 cross-references occur inside the counted spans; they are listed in external\_references as navigation pointers only.\par\medskip
\noindent\textit{Allowed standard prerequisites (editorial).}\begin{itemize}
\item A scheme carrying an ample invertible sheaf is quasi-compact, and the existence of a norm of degree d implies that the morphism is surjective. Cited at proof lines 3746-3749. Gives quasi-compactness of the base, which is the first half of the ampleness criterion being applied.
\item Ampleness tested by affine caps: it suffices to find, at every point, a section of a positive tensor power not vanishing at that point whose nonvanishing locus is affine, and the existence of a section of a high power of an ample sheaf whose nonvanishing locus separates a finite set from a prescribed affine open. Cited at proof lines 3752-3761. Produces the section downstairs whose nonvanishing locus lies inside a chosen affine open and whose behaviour at the norm is controlled.
\item The norm of a section: for a finite morphism with a degree d norm the norm of a section of a power of an invertible sheaf is a section of the same power of the norm, whose vanishing behaviour is controlled by the section's vanishing behaviour. Cited at proof lines 3762-3764. Identifies the candidate sheaf on the base as ample and gives the nonvanishing statement.
\item The nonvanishing locus of a section of a sheaf restricted to an affine open is affine. Cited at proof lines 3765-3767. Concludes that the cap of the chosen section is affine, completing the ampleness criterion.
\end{itemize}\medskip
\ifdefined\hypertarget\hypertarget{sourcepointer-1}{}\fi
\noindent\textit{Source reference.} Original lemma 0BCZ: focused-seeds/AGP/sources/stacks-snapshot-20261011/tags.txt, source lines 14297--14297. Same-chapter reference used in the retained author text at source line(s) 3768. It is supplied only as a navigation pointer so that the retained text compiles; the counted item delivers the author's own statement and the author's complete proof, and no part of the referenced item is copied in or substituted for the counted argument.\par\medskip
\ifdefined\hypertarget\hypertarget{sourcepointer-2}{}\fi
\noindent\textit{Source reference.} Original definition 01PS: focused-seeds/AGP/sources/stacks-snapshot-20261011/tags.txt, source lines 2241--2241. Same-chapter reference used in the retained author text at source line(s) 3751. It is supplied only as a navigation pointer so that the retained text compiles; the counted item delivers the author's own statement and the author's complete proof, and no part of the referenced item is copied in or substituted for the counted argument.\par\medskip
\ifdefined\hypertarget\hypertarget{sourcepointer-3}{}\fi
\noindent\textit{Source reference.} Original lemma 01PV: focused-seeds/AGP/sources/stacks-snapshot-20261011/tags.txt, source lines 2244--2244. Same-chapter reference used in the retained author text at source line(s) 3769. It is supplied only as a navigation pointer so that the retained text compiles; the counted item delivers the author's own statement and the author's complete proof, and no part of the referenced item is copied in or substituted for the counted argument.\par\medskip
\ifdefined\hypertarget\hypertarget{sourcepointer-4}{}\fi
\noindent\textit{Source reference.} Original lemma 09NV: focused-seeds/AGP/sources/stacks-snapshot-20261011/tags.txt, source lines 12228--12228. Same-chapter reference used in the retained author text at source line(s) 3254, 3765. It is supplied only as a navigation pointer so that the retained text compiles; the counted item delivers the author's own statement and the author's complete proof, and no part of the referenced item is copied in or substituted for the counted argument.\par\medskip
\begin{lemma}
\label{lemma-norm-ample}
Let $\pi : X \to Y$ be a finite morphism of schemes. Assume $X$ has
an ample invertible sheaf and there exists a norm of degree $d$
for $\pi$. Then $Y$ has an ample invertible sheaf.
\end{lemma}

\begin{proof}
Let $\mathcal{L}$ be the ample invertible sheaf on $X$ given to us
by assumption. We will prove that $\mathcal{N} = \text{Norm}_\pi(\mathcal{L})$
is ample on $Y$.

\medskip\noindent
Since $X$ is quasi-compact (Properties, Definition
\ref{properties-definition-ample}) and $X \to Y$ surjective
(by the existence of $\text{Norm}_\pi$)
we see that $Y$ is quasi-compact.
Let $y \in Y$ be a point. To finish the proof
we will show that there exists a section $t$ of some positive tensor
power of $\mathcal{N}$ which does not vanish at $y$ such that $Y_t$
is affine. To do this, choose an affine open neighbourhood $V \subset Y$
of $y$. Choose $n \gg 0$ and a section
$s \in \Gamma(X, \mathcal{L}^{\otimes n})$
such that
$$
\pi^{-1}(\{y\}) \subset X_s \subset \pi^{-1}V
$$
by
Properties, Lemma \ref{properties-lemma-ample-finite-set-in-principal-affine}.
Then $t = \text{Norm}_\pi(s)$ is a section of $\mathcal{N}^{\otimes n}$
which does not vanish at $x$ and with $Y_t \subset V$, see
Lemma \ref{lemma-norm-map-invertible}. Then $Y_t$
is affine by Properties, Lemma \ref{properties-lemma-affine-cap-s-open}.
\end{proof}

\end{document}
